% =============================================================================
\section{Conclusion}
\label{sec:conclusion}
% =============================================================================

We introduced three bidirectional strategies for LLM-based setwise document ranking that challenge the convention of extracting only ``best'' information from each LLM comparison. Our experiments with nine LLMs across three architecture families yield three findings.

First, \textbf{the DualEnd family is the strongest overall}. \dualend{}-Cocktail achieves the highest average NDCG@10 across all 9 models on both DL 2019 (.7134) and DL 2020 (.6791), and DualEnd methods (Cocktail or Selection) win in 14 of 18 model--dataset configurations. The per-query significance picture is more conservative: only one of those gains remains Bonferroni-significant. By asking the LLM for both the best and worst document per call, dual-end methods narrow the ranking from both ends, and an unexpected benefit is \emph{implicit debiasing}: dual-end prompting produces more uniform position selection patterns than single-direction methods, partially mitigating the strong recency bias observed in standalone best- and worst-selection.

Second, \textbf{bottom-up ranking is substantially weaker than top-down}, refuting the hypothesis that LLMs are better at identifying irrelevant documents. The gap is particularly severe for smaller models (Flan-T5-Large: .289 vs.\ .654 NDCG@10) and narrows but does not vanish at scale. Bonferroni-significant losses appear in 6 of 18 configurations. Position bias analysis reveals a likely mechanism: models exhibit extreme recency bias when selecting ``worst'' (position D selected 40--63\% of the time), effectively defaulting to the last passage.

Third, \textbf{bidirectional ensemble does not outperform the top-down baseline}, indicating that simply fusing top-down and bottom-up rankings cannot compensate for the weakness of the bottom-up component. Even at $\alpha=0.9$ (90\% TopDown weight), fusion underperforms pure TopDown, and BiDir produces 3 Bonferroni-significant losses. The intuitive assumption that directional fusion would parallel the success of permutation voting~\citep{tang2024found} does not hold because the directions are not equally reliable.

The strongest scientific contribution is therefore one of \emph{directional asymmetry}: worst-selection alone is not the dual of best-selection, and the gap between standalone worst-selection and joint best-worst elicitation is mechanistically meaningful. The strongest practical contribution is the empirical map that says which sorting algorithm and which selection direction works for which model family, and the quality--cost Pareto frontier that pinpoints the missing intermediate point between \topdown{}-Bubble and \dualend{}-Cocktail.

Our findings open several avenues for future work: (1) full empirical evaluation of the Selective, bias-aware, and same-call regularized DualEnd refinement variants, all of which are already implemented and instrumented in our fork; (2) extending dual-end to extract a \emph{full ordering} of the candidate set, effectively bridging setwise and listwise paradigms; (3) investigating whether reasoning-enhanced rerankers~\citep{zhuang2025rankr1} benefit from dual-end prompting; (4) benchmarking the new Qwen/Qwen3.5 likelihood follow-up against the generation results to quantify the cost saved and the quality lost; and (5) understanding why worst-selection is harder for LLMs and whether this can be mitigated through chain-of-thought reasoning.
